% ============================================================
\appendix
\chapter{Configuration Reference}
\label{app:configuration}
% ============================================================

\section{PIDNet-S Configuration}

The portable YAML below is the configuration used for verified evaluation. It
uses repository-relative paths; the training checkpoint contains the original
full configuration.

\begin{lstlisting}
seed: 1337
raw_root: CAT
processed_root: data_processed_blue_green
reports_dir: reports
configs_dir: configs
outputs_dir: outputs
dataset_name: mixed
experiment_name: controlled_ablation_augment_off

preprocessing:
  val_fraction: 0.2
  traversable_palette_names: [blue, green]
  mapping_filename: discovered_mapping_blue_green.yaml

training:
  model_name: pidnet-s
  input_size: [640, 384]
  batch_size: 4
  epochs: 15
  learning_rate: 0.00025
  min_learning_rate: 0.00001
  weight_decay: 0.0001
  mixed_precision: bf16
  num_workers: 0
  augment: false
  untraversable_weight_bias: 2.5
  traversable_weight_bias: 0.75

postprocessing:
  traversable_threshold: 0.5
\end{lstlisting}

\section{Reconstructed ROD Configuration}

\begin{lstlisting}
seed: 1337
raw_root: CAT
processed_root: data_processed_blue_green
reports_dir: reports
configs_dir: configs
outputs_dir: outputs
dataset_name: mixed
experiment_name: mixed_binary_traversability_rod_vits_paper

preprocessing:
  val_fraction: 0.2
  traversable_palette_names: [blue, green]
  mapping_filename: discovered_mapping_blue_green.yaml

training:
  model_name: rod-vits
  rod_encoder_checkpoint: weights/efficient_sam_vits.pt
  rod_encoder_image_size: 1024
  input_size: [640, 384]
  batch_size: 8
  val_batch_size: 4
  epochs: 15
  learning_rate: 0.001
  weight_decay: 0.01
  lr_scheduler: poly
  poly_power: 0.9
  mixed_precision: bf16
  num_workers: 4
  augment: false
  class_weights: [1.0, 1.0]
  ce_loss_weight: 1.0
  dice_loss_weight: 0.0

postprocessing:
  traversable_threshold: 0.5
\end{lstlisting}

\section{Common Fixed-Budget Architecture-Comparison Template}

The model name, experiment name, seed, processed root, and mapping filename
change by row and policy. The remaining controlled settings follow this
template:

\begin{lstlisting}
training:
  model_name: <architecture-specific name>
  encoder_weights: <imagenet | null | model-specific>
  input_size: [640, 384]
  batch_size: 4
  epochs: 15
  learning_rate: 0.00025
  min_learning_rate: 0.00001
  weight_decay: 0.0001
  mixed_precision: bf16
  num_workers: 4
  augment: true
  untraversable_weight_bias: 2.5
  traversable_weight_bias: 0.75
postprocessing:
  traversable_threshold: 0.5
\end{lstlisting}

Table~\ref{tab:appendix_model_names} maps each human-readable architecture
name to the exact configuration string consumed by the model factory.

\begin{table}[H]
\centering
\caption{Configuration model names used to construct each architecture-comparison
system.}
\label{tab:appendix_model_names}
\small
\begin{tabularx}{\textwidth}{@{}lX@{}}
\toprule
\textbf{System} & \textbf{\texttt{training.model\_name}} \\
\midrule
PIDNet-S & \texttt{pidnet-s} \\
ROD ViT-S & \texttt{rod-vits} \\
FPN/MobileNetV2 & \texttt{smp:FPN:mobilenet\_v2} \\
FPN/EfficientNet-B0 & \texttt{smp:FPN:efficientnet-b0} \\
U-Net/MobileNetV2 & \texttt{smp:Unet:mobilenet\_v2} \\
U-Net/EfficientNet-B0 & \texttt{smp:Unet:efficientnet-b0} \\
SegFormer-B0 & \texttt{smp:Segformer:mit\_b0} \\
DDRNet-23-Slim & \texttt{ddrnet-23-slim} \\
BiSeNetV2 & \texttt{bisenetv2} \\
\bottomrule
\end{tabularx}
\end{table}

\section{Convergence-Aware Follow-Up Template}

The convergence suite retains the model-specific fields and common settings
above, but replaces the fixed horizon with the following schedule and stopping
rule:

\begin{lstlisting}
training:
  epochs: 300                 # ceiling, not claimed optimum
  lr_scheduler: cosine
  cosine_decay_epochs: 60     # then hold min_learning_rate
  save_last_checkpoint: false
  save_optimizer_state: false
  early_stopping:
    enabled: true
    metric: mIoU
    mode: max
    patience: 25
    min_epochs: 60
    min_delta: 0.0001
postprocessing:
  traversable_threshold: 0.5
\end{lstlisting}

The complete 27 generated YAML files are stored under
\path{configs/generated/convergence_blue_green_e300_c60_m60_p25/blue_green/}.
They should be used instead of reconstructing a run from this shortened
listing.

% ============================================================
\chapter{Reference Evaluation Record}
\label{app:expected}
% ============================================================

The following values are sufficient to check an independent checkpoint
evaluation. The confusion matrix is ordered by ground-truth row and predicted
column.

\begin{lstlisting}
split: test
checkpoint: outputs/controlled_ablation_augment_off/checkpoints/best.pt
threshold: 0.5
confusion_matrix:
  - [73151244, 3300886]
  - [4027996, 53213314]
mIoU: 0.8939406948414585
iou_untraversable: 0.9089355053942137
iou_traversable: 0.8789458842887033
precision_traversable: 0.9415919184912819
recall_traversable: 0.9296313099752609
f1_traversable: 0.9355733889285891
false_safe_rate: 0.04317585396247299
false_block_rate: 0.07036869002473913
\end{lstlisting}

The TensorRT PIDNet result is expected to differ slightly:

\begin{lstlisting}
confusion_matrix:
  - [73155734, 3296396]
  - [4045047, 53196263]
mIoU: 0.8937640836927854
false_safe_rate: 0.04311712440189698
false_block_rate: 0.07066656930108693
\end{lstlisting}

The reconstructed ROD TensorRT full-test record is:

\begin{lstlisting}
confusion_matrix:
  - [73508109, 2944021]
  - [2265250, 54976060]
mIoU: 0.9236346417584864
false_safe_rate: 0.03850803110390776
false_block_rate: 0.03957369249585658
\end{lstlisting}

% ============================================================
\chapter{Reproduction Commands}
\label{app:reproduction}
% ============================================================

\noindent\textbf{Reproducibility repository:}\\
\url{https://github.com/hax2/cat-pidnet-thesis-reproducibility}. The repository
contains the code, configurations, checkpoints and evidence records, and the
instructions needed for reproduction. Run the commands in this appendix from
its root. The convergence evidence was generated from source commit
\texttt{865cda90}; later thesis-only edits do not alter the recorded training
artifacts or metrics.

\section{Audit and Preprocessing}

From the repository root, the raw dataset can be audited and the main binary
dataset regenerated with:

\begin{lstlisting}
PYTHONPATH=src python -m ttfm.cli \
  --config configs/frozen_pidnet_cat.yaml audit

PYTHONPATH=src python -m ttfm.cli \
  --config configs/frozen_pidnet_cat.yaml preprocess
\end{lstlisting}

Preprocessing replaces the configured processed root, so the retained manifest
identity should be checked before regenerating evidence.

\section{PIDNet-S Evaluation}

\begin{lstlisting}
PYTHONPATH=src python -m ttfm.cli \
  --config configs/frozen_pidnet_cat.yaml \
  eval --split test
\end{lstlisting}

This command loads \texttt{best.pt}; it does not update weights. After
processing the complete test manifest, it writes \path{test_metrics.json}.

\section{Ranked Qualitative Review}

\begin{lstlisting}
PYTHONPATH=src python -m ttfm.cli \
  --config configs/frozen_pidnet_cat.yaml \
  review --split test --top-k 10
\end{lstlisting}

The ALICE panel is regenerated separately:

\begin{lstlisting}
PYTHONPATH=src python scripts/alice_preview.py \
  --config configs/frozen_pidnet_cat.yaml \
  --input-dir alice \
  --output-dir outputs/alice_pidnet_frozen_preview \
  --sample-count 12 \
  --checkpoint best.pt
\end{lstlisting}

\section{TensorRT Benchmarks}

\begin{lstlisting}[basicstyle=\ttfamily\tiny]
PYTHONPATH=src python scripts/edge_trt_benchmark.py \
  --config configs/ablations/controlled_ablation_augment_off.yaml \
  --checkpoint outputs/controlled_ablation_augment_off/checkpoints/best.pt \
  --input-dir data_processed_blue_green/test/images \
  --precision fp16 --limit 64 --warmup 10 --repeats 5 \
  --accuracy-limit 0 \
  --output reports/edge_pidnet_rtx5060_trt.json
\end{lstlisting}

\begin{lstlisting}[basicstyle=\ttfamily\tiny]
PYTHONPATH=src python scripts/edge_trt_benchmark.py \
  --config rod_suite_results/configs/rod_vits_cat_paper.yaml \
  --checkpoint rod_suite_results/outputs/mixed_binary_traversability_rod_vits_paper/checkpoints/best.pt \
  --input-dir data_processed_blue_green/test/images \
  --precision fp16 --limit 64 --warmup 10 --repeats 5 \
  --accuracy-limit 0 \
  --output reports/edge_rod_rtx5060_trt.json
\end{lstlisting}

Both commands require a CUDA and TensorRT environment. Engine construction is
recorded but excluded from inference timing. Exact latency is
hardware- and runtime-dependent even when accuracy is reproducible.

\section{Architecture Ledgers}

The readable aggregate reports are:
\begin{lstlisting}
reports/blue_only_verified_summary.md
reports/blue_green_verified_summary.md
reports/convergence_blue_green_verified_summary.md
\end{lstlisting}
Their JSON ledgers contain per-seed configurations, metrics, confusion counts,
artifact paths, and integrity checks. The convergence ledger additionally
contains validation-selected epochs and server-recorded checkpoint identities.
A selected run can be reevaluated after retrieving its checkpoint with:
\begin{lstlisting}
PYTHONPATH=src python -m ttfm.cli \
  --config <exact-config.yaml> eval --split test
\end{lstlisting}

% ============================================================
\chapter{Artifact Identities}
\label{app:hashes}
% ============================================================

These SHA-256 values identify the evidence package at the time of verification.
A different digest denotes a different artifact, even if its filename is
unchanged.
Table~\ref{tab:appendix_hashes} records the retained PIDNet evidence chain.

\begin{longtable}{@{}p{0.46\textwidth}p{0.48\textwidth}@{}}
\caption{SHA-256 identities for the PIDNet evidence chain.}\label{tab:appendix_hashes}\\
\toprule
\textbf{Artifact} & \textbf{SHA-256} \\
\midrule
\endfirsthead
\toprule
\textbf{Artifact} & \textbf{SHA-256} \\
\midrule
\endhead
\path{configs/frozen_pidnet_cat.yaml} &
\texttt{\seqsplit{07b07bbb1bb9e6f0586df6598d2cf982f8af845fb3e16785810615c80ebdaa74}} \\
\path{data_processed_blue_green/manifest.json} &
\texttt{\seqsplit{198fef93fc45568e1dae49ae8780906fa6ef479548214d20244cd12e51201fa2}} \\
\path{outputs/.../checkpoints/best.pt} &
\texttt{\seqsplit{b9862fe49713c0622d760b94aa3fc055fdaf6549a67a4ee3492004d9c591f471}} \\
\path{src/ttfm/eval.py} &
\texttt{\seqsplit{f83a6da0902bec7226a0a23112c97fdf549ce493642a002b535285c7d18fdd78}} \\
\path{src/ttfm/metrics.py} &
\texttt{\seqsplit{84ee42f674b97ae7f14b04d68e09c9c899c89c8d528f178897ec67610d8662fc}} \\
\path{src/ttfm/review.py} &
\texttt{\seqsplit{678820d64404d8a3f5ed2b9c59a54ac63596c67acb4082df53bc1a78604c73ee}} \\
\path{outputs/.../test_metrics.json} &
\texttt{\seqsplit{f72f173c5c64bc2b4e5ff11bffc9184962db7f33fe2cc82136c661282bdd2263}} \\
\path{outputs/.../test_review/summary.json} &
\texttt{\seqsplit{bc0b0237acd555c4361540afa1ab0ad0b445734662877fb52fe77a37ebaca0b2}} \\
\path{reports/cpu_pidnet_s_ryzen5500_torch213_idle.json} &
\texttt{\seqsplit{4db19ce68688fcdca4a44924dfa64670d9045474448703feff3302c459f53f5b}} \\
\path{configs/ablations/controlled_ablation_augment_off.yaml} &
\texttt{\seqsplit{62869964e3fd30b3cba59d28f1799c6add73002c1a39af4dfb32641626d27928}} \\
\path{scripts/edge_trt_benchmark.py} &
\texttt{\seqsplit{0a1acc1a55659384208cbf067f3c90c7c5a304b47749bada65e7be6cc624f2a3}} \\
\path{reports/.trt_scratch/pidnet-s_fp16.onnx} &
\texttt{\seqsplit{15ac7c3ea859e8776b073aaafca76a5507f501cd62be992efc54d08276d43d91}} \\
\path{reports/.trt_scratch/pidnet-s_fp16.engine} &
\texttt{\seqsplit{2458b84daa0133d0de7fea5f2c20fd1ecb3f335eaa00e81bd6f00c8f3f52da85}} \\
\path{reports/edge_pidnet_rtx5060_trt.json} &
\texttt{\seqsplit{5eb085823c892bc40d34d4997df5745206246ee8eba14cf906a7619d43107255}} \\
\bottomrule
\end{longtable}

% ============================================================
\chapter{Implementation Map}
\label{app:implementation}
% ============================================================

Table~\ref{tab:appendix_implementation} maps the main repository paths to
their role in creating or checking the experimental record.

\begin{longtable}{@{}p{0.34\textwidth}p{0.60\textwidth}@{}}
\caption{Principal repository files and their role in the experimental record.}\label{tab:appendix_implementation}\\
\toprule
\textbf{File or directory} & \textbf{Responsibility} \\
\midrule
\endfirsthead
\toprule
\textbf{File or directory} & \textbf{Responsibility} \\
\midrule
\endhead
\path{src/ttfm/audit.py} & Discover file counts, IDs, palette colours, dimensions, pairing, and corruption. \\
\path{src/ttfm/preprocess.py} & Apply the binary mapping, deterministic hash split, and write manifests and binary masks. \\
\path{src/ttfm/data.py} & Load manifest samples, resize RGB/masks, apply deterministic augmentation, scale tensors, and compute train-only class weights. \\
\path{src/ttfm/model.py} & Construct PIDNet, ROD, DDRNet, BiSeNet, and SMP systems behind one full-resolution logits interface. \\
\path{src/ttfm/vendor_pidnet/} & Vendored PIDNet implementation used by the baseline wrapper. \\
\path{src/ttfm/rod.py} & EfficientSAM ViT-S encoder, frozen-weight loading, intermediate feature extraction, ROD decoder, and output resizing. \\
\path{src/ttfm/train.py} & Losses, AdamW, cosine/polynomial schedules, mixed precision, epoch loops, validation, and best-checkpoint selection. \\
\path{src/ttfm/eval.py} & Read-only checkpoint evaluation and JSON metric output. \\
\path{src/ttfm/metrics.py} & Softmax threshold, confusion accumulation, mIoU, class IoUs, precision, recall, F1, FSR, and FBR. \\
\path{src/ttfm/review.py} & Per-image asymmetric-error and boundary scoring plus ranked qualitative panels. \\
\path{scripts/run_rod_suite.py} & Controlled PIDNet--ROD seeds, reconstructed training runs, common GPU timing, and aggregate reports. \\
\path{scripts/run_realtime_comparison_suite.py} & Nine-model blue-only repeat-seed comparison. \\
\path{scripts/run_realtime_blue_green_suite.sh} & Corresponding blue+green architecture suite. \\
\path{scripts/run_convergence_suite.py} & Generate, resume, evaluate, and aggregate the 27-run convergence-aware blue+green suite. \\
\path{scripts/build_convergence_evidence.py} & Audit convergence artifacts and regenerate the verified ledger, summary, and figures. \\
\path{scripts/gpu_inference_benchmark.py} & Eager/CPU forward benchmark with warm-up and repeat controls. \\
\path{reports/cpu_*_ryzen5500*.json} & Retained eager and compiled Ryzen benchmark records, including thread count, warm-up, repeats, sample count, and latency percentiles. \\
\path{scripts/edge_trt_benchmark.py} & ONNX export, TensorRT engine construction, scoped latency, and full-test engine accuracy. \\
\path{reports/*verified_ledger.json} & Per-seed metric, configuration, checkpoint, benchmark, and integrity records. \\
\bottomrule
\end{longtable}

% ============================================================
\chapter{Development Time and Work Plan}
\label{app:development_time}
% ============================================================

The project was developed iteratively rather than from a fixed recipe. Early
work compared several possible formulations of off-road perception before the
binary CaT task and the final architecture study were selected. Table
\ref{tab:development_time} reconstructs the effort devoted to each principal
step. The values are approximate because a formal timesheet was not maintained;
they represent about 22 working days, or roughly 180 hours, of research,
implementation, experimentation, and writing. Reading and revision continued
across several phases, so the entries should be interpreted as an effort
breakdown rather than as mutually exclusive calendar blocks.

\begingroup
\small
\begin{longtable}{@{}p{0.25\textwidth}p{0.17\textwidth}p{0.50\textwidth}@{}}
\caption[Development-time breakdown.]{Approximate development time devoted to
the principal stages of the thesis. Automated computation is distinguished
from hands-on research work where it materially affects the estimate.}
\label{tab:development_time} \\
\toprule
\textbf{Stage} & \textbf{Estimated effort} & \textbf{Work performed and reason for the effort} \\
\midrule
\endfirsthead
\toprule
\textbf{Stage} & \textbf{Estimated effort} & \textbf{Work performed and reason for the effort} \\
\midrule
\endhead
\midrule
\multicolumn{3}{r}{\textit{Continued on the next page}} \\
\endfoot
\bottomrule
\endlastfoot
Problem framing, literature review, and alternative approaches &
4 days (about 30 h) &
Reviewed monocular, geometric, multimodal, and segmentation-based approaches;
compared possible thesis directions; traced relevant CaT and ROD literature;
and contacted an original author to query apparent inconsistencies in the
published experimental description. \\

Dataset search and selection &
3 days (about 23 h) &
Compared candidate off-road datasets by annotation type, licence, scene
diversity, task compatibility, and reproducibility before selecting CaT. \\

CaT audit and dataset preparation &
2 days (about 16 h) &
Audited RGB--mask--integer-map correspondence, identified error-prone filename
pairing assumptions, reconstructed the palette mapping, defined the two binary
policies, and produced deterministic development splits. \\

PIDNet-S baseline development &
2 days (about 16 h) &
Adapted PIDNet-S to two-class output, implemented the common training and
evaluation path, checked tensor shapes and losses, and obtained a reproducible
working baseline. \\

Controlled ablations and checkpoint audit &
2 days (about 16 h) &
Ran the controlled variants, checked validation-based checkpoint selection,
recomputed strict test metrics, and generated qualitative error reviews. \\

ROD reconstruction and comparison &
1 day (about 8 h) &
Reconstructed the published frozen-encoder recipe, adapted it to the CaT
policies, and established the matched PIDNet--ROD comparison. \\

Architecture extension and reproducibility evidence &
1.5 days (about 12 h) &
Integrated seven further compact systems, standardised their interfaces and
configurations, executed the multi-seed matrices, and assembled traceable
result ledgers. \\

Hardware evaluation &
1 day (about 8 h) &
Defined timing scopes and measured CPU, H100, and TensorRT execution while
checking that exported models preserved their source predictions. \\

Convergence-aware follow-up &
1 full day (about 13 h elapsed) &
Prepared, launched, monitored, and verified 27 long-horizon runs. Early
stopping reduced the 300-epoch ceiling to validation-selected checkpoints from
epochs 50--141. \\

Writing, figure production, and final revision &
4.5 days (about 36 h) &
Organised the argument, documented equations and protocols, produced and
revised figures and tables, incorporated the convergence evidence, and checked
language, citations, reproducibility, and UC3M presentation requirements. \\
\midrule
\textbf{Total} & \textbf{about 22 days (about 180 h)} &
Approximate hands-on and supervised experimental effort for the complete
thesis. \\
\end{longtable}
\endgroup
